\chapter{Challenge 4: Parkmanöver}
\label{ch:Challenge_4}

% Einleitung zur Challenge 4
Im Rahmen dieser Aufgabe wurde der Duckiebot mit der Anforderung konfrontiert, eine freie Parklücke zu identifizieren, 
rückwärts einzuparken, fünf Sekunden im Ruhezustand zu verweilen und anschließend die Fahrt wieder aufzunehmen. Die 
Herausforderung bestand darin, eine geeignete Fläche innerhalb des Verkehrsraums von Duckietown zu detektieren, die 
weder durch Enten noch durch andere Duckiebots belegt ist. Basierend auf den zuvor entwickelten Modulen für Spurverfolgung 
und Kreuzungshandhabung musste ein Steuerungsmechanismus zur Durchführung des Parkvorgangs implementiert werden. 
Ziel war es, die Fähigkeit zur temporären Deaktivierung des Fahrbetriebs zu simulieren, wie sie beispielsweise in 
Ladesituationen oder bei planmäßigen Stopps erforderlich ist. 